\subsection{Fault-Tolerant Quantum Gates}

So far, we have studied quantum computation in an idealized setting, where quantum gates are represented by perfect unitary matrices acting on perfectly isolated qubits. In practice, however, real quantum hardware is extremely fragile. Physical qubits are affected by noise, decoherence, imperfect control operations, and measurement errors, all of which can corrupt the quantum information stored in the system.

\highspace
To perform reliable quantum computation, it is therefore necessary to protect quantum information against errors. This is the goal of \textbf{quantum error correction}, where a single logical qubit is encoded into many physical qubits in such a way that errors can be detected and corrected without destroying the quantum state.

\highspace
However, protecting quantum information is not sufficient by itself. We must also be able to apply quantum gates directly on encoded logical qubits while preserving the fault-tolerant structure of the error-correcting code. This introduces the notion of \textbf{fault-tolerant quantum gates}, namely gates designed so that local physical errors do not spread uncontrollably through the computation.

\highspace
In this context, Clifford gates play a particularly important role because many of them admit simple fault-tolerant implementations, often through \textbf{transversal operations}. Non-Clifford gates, on the other hand, are significantly more difficult to implement fault-tolerantly, even though they are essential for universal quantum computation. This creates one of the central challenges of modern quantum computing: the gates required for computational universality are also the most expensive and difficult to protect against errors.

\highspace
In the following sections, we introduce the distinction between physical and logical qubits, the concept of transversal gates, and the reasons why non-Clifford gates such as the $T$ gate represent the main obstacle in fault-tolerant quantum computation.